\scanpage{13}
\mspage{10}

plus \unsure{exactement|exactemt} sur $W\times I$, et si \deleted{\unsure{c'est}} $\alpha$

\[
\alpha:W\times I\longrightarrow\unsure{\mathbb{\Delta}|\Delta}_{1}\times\mathbb{\Delta}_{1}\times I
\]

est l'inclusion canonique, on définira le carré ``cartésien'' comme \inserted{l'image par} $\alpha_{*}$ de ce coéfficient. Cela nous ramène donc au problème de ``globaliser'' \inserted{d'abord} le diagramme partiel

\diagram{\begin{tikzcd}[ampersand replacement=\&]
F\arrow[dr] \& \\ \& j_{*}j^{*}F \\ e\arrow[ur] \& {.}
\end{tikzcd}}

Et de $=$ pour le carré (1), sauf qu'on est embarrassé si on choisit de le regarder \unsure{surtout|seulemt} comme cartésien (ce qui suggère de procéder comme tantôt) ou comme \uline{co}cartésien (ce qui suggérerait de procéder de façon duale). Mais les deux coeff. qu'on trouve par ces deux procédés, \deleted{devraient être \unsure{canoniques}} \inserted{en partant soit de}

\diagram{\begin{tikzcd}[ampersand replacement=\&]
F\arrow[dr,very thick] \& \\ \& i_{*}i^{*}F \\ e\arrow[ur] \& {,}
\end{tikzcd}}

soit de

\diagram{\begin{tikzcd}[ampersand replacement=\&]
\& F \\ j_{!}j^{*}F\arrow[ur]\arrow[dr] \& \\ \& e,
\end{tikzcd}}

devraient être \uline{canoniquement} isomorphes.

\deleted{Il y a lieu d'ailleurs de \unsure{privilégier|préférer} le premier, qui ne fait pas appel à un biadjoint, et d'en déduire le deuxième (via \unsure{prolongmt} du coeff. sur $W\times I$ à un coeff sur $\mathbb{\Delta}_{1}\times\mathbb{\Delta}_{1}\times I$), d'où aussi une \unsure{détermination|description} \unsure{du|de} $\unsure{j_{!}|j_{*}}j^{*}F$}

% archive sheet number: 12
